\subsection{MAURER-CARTAN FORMULAE}

Let X be a manifold of class \(C^{r}\) , of finite dimension if K is of characteristic \textgreater0, and let L be a complete normable Lie algebra. Let \(\alpha\) be a differential form of degree 1 on X with values in L of class \(C^{r-1}\) . Let \(x \in X\) . The mapping

\[
(u _ {1}, u _ {2}) \mapsto [ \alpha_ {x} (u _ {1}), \alpha_ {x} (u _ {2}) ]
\]

of \(\mathrm{T}_x(\mathbf{X})\times \mathrm{T}_x(\mathbf{X})\) into L is a continuous alternating bilinear form on \(\mathrm{T}_x(\mathbf{X})\) with values in L. We shall denote it by \([\alpha ]_x^2\) so that \([\alpha ]^2\) is a differential form of degree 2 on X with values in L. Identifying an open neighbourhood of \(x\) in X with an open subset of a Banach space, we see immediately that \([\alpha ]^2\) is of class \(\mathrm{C}^{r - 1}\) . If \(\mathbf{X}'\) is a manifold of class \(\mathrm{C}^r\) and \(f\colon \mathbf{X}'\to \mathbf{X}\) is a morphism, then

\[
[ f ^ {*} (\alpha) ] ^ {2} = f ^ {*} ([ \alpha ] ^ {2}).\tag{28}
\]

Let \(\alpha, \beta\) be two differential forms of degree 1 on X with values in L of class \(C^{r-1}\) . The exterior product \(\alpha \wedge \beta\) of \(\alpha\) and \(\beta\) (Differentiable and Analytic Manifolds, R, 7.8.2) is a differential form of degree 2 on X with values in L of class \(C^{r-1}\) ; we have

\[
(\alpha \wedge \beta) _ {x} (u _ {1}, u _ {2}) = [ \alpha_ {x} (u _ {1}), \beta_ {x} (u _ {2}) ] - [ \alpha_ {x} (u _ {2}), \beta_ {x} (u _ {1}) ]\tag{29}
\]

for \(u_{1}, u_{2}\) in \(\mathrm{T}_x(\mathbf{X})\) . It is immediate that

\begin{enumerate}
\def\labelenumi{(\arabic{enumi})}
\setcounter{enumi}{29}
\tightlist
\item
\end{enumerate}

\[
[ \alpha + \beta ] ^ {2} = [ \alpha ] ^ {2} + [ \beta ] ^ {2} + \alpha \wedge \beta\tag{31}
\]

\[
\alpha \wedge \alpha = 2 [ \alpha ] ^ {2}.
\]

PROPOSITION 51. Let G be a Lie group, of finite dimension if K is of characteristic \textgreater0, and let \(a_{1},\ldots,a_{p}\) be elements of L(G), F a complete normable space and \(\alpha\) a differential form of degree \(p-1\) on G with values in F. If \(\alpha\) is left invariant, then

\[
(d \alpha) _ {e} (a _ {1}, \dots , a _ {p}) = \sum_ {i <   j} (- 1) ^ {i + j} \alpha_ {e} ([ a _ {i}, a _ {j} ], a _ {1}, \dots , a _ {i - 1}, a _ {i + 1}, \dots , a _ {j - 1},
\]

\[
\left. a _ {j + 1}, \dots , a _ {p}\right).
\]

If \(\alpha\) is right invariant, then

\[
\begin{array}{l} (d \alpha) _ {e} (a _ {1}, \dots , a _ {p}) \\ = - \sum_ {i <   j} (- 1) ^ {i + j} \alpha_ {e} ([ a _ {i}, a _ {j} ], a _ {1}, \dots , a _ {i - 1}, a _ {i + 1}, \dots , a _ {j - 1}, a _ {j + 1}, \dots , a _ {p}). \end{array}
\]

Suppose that \(\alpha\) is left invariant. By Differentiable and Analytic Manifolds, R, 8.5.7, then

\[
\begin{array}{l} (d \alpha) (L _ {a _ {1}}, \ldots , L _ {a _ {p}}) = \sum_ {i} (- 1) ^ {i - 1} L _ {a _ {i}} \alpha (L _ {a _ {1}}, \ldots , L _ {a _ {i - 1}}, L _ {a _ {i + 1}}, \ldots , L _ {a _ {p}}) \\ + \sum_ {i <   j} (- 1) ^ {i + j} \alpha ([ L _ {a _ {i}}, L _ {a _ {j}} ], L _ {a _ {1}}, \ldots , L _ {a _ {i - 1}}, L _ {a _ {i + 1}}, \ldots , L _ {a _ {j - 1}}, L _ {a _ {j + 1}}, \ldots , L _ {a _ {p}}). \end{array}
\]

But the functions \(\alpha (\mathrm{L}_{a_1},\dots ,\mathrm{L}_{a_{i - 1}},\mathrm{L}_{a_{i + 1}},\dots ,\mathrm{L}_{a_p})\) on G are left invariant and therefore constant. Hence

\[
\mathrm{L} _ {a _ {i}} \alpha (\mathrm{L} _ {a _ {1}}, \dots , \mathrm{L} _ {a _ {i - 1}}, \mathrm{L} _ {a _ {i + 1}}, \dots , \mathrm{L} _ {a _ {p}}) = 0.
\]

Moreover, \([L_{a_{i}}, L_{a_{j}}] = L_{\{a_{i}, a_{j}\}}\) (Proposition 23), whence the first formula of Proposition 51. The second can be established analogously, this time using the relation \([R_{a_{i}}, R_{a_{j}}] = -R_{\{a_{i}, a_{j}\}}\) .

COROLLARY 1. Let G be a Lie group, of finite dimension if K is of characteristic \textgreater0, and ω and ω' the left and right canonical differential forms of G. Then

\[
d \omega + [ \omega ] ^ {2} = 0, \quad d \omega^ {\prime} - [ \omega^ {\prime} ] ^ {2} = 0.
\]

By Proposition 51,

\[
\begin{array}{r l} (d \omega) _ {e} (a _ {1}, a _ {2}) & = - \omega_ {e} ([ a _ {1}, a _ {2} ]) = - [ a _ {1}, a _ {2} ] = - [ \omega_ {e} (a _ {1}), \omega_ {e} (a _ {2}) ] \\ & = - [ \omega ] _ {e} ^ {2} (a _ {1}, a _ {2}) \end{array}
\]

whence the first formula. The second can be established analogously.

COROLLARY 2. Suppose that G is finite-dimensional. Let \((e_1, \ldots, e_n)\) be a basis of L(G), \((e_1^*, \ldots, e_n^*)\) the dual basis, \((c_{ijk})\) the constants of structure of L(G) relative to the basis \((e_1, \ldots, e_n)\) and \(\omega_i\) (resp. \(\omega_i^{\prime})\) the left (resp. right) invariant differential form on G with values in K such that \((\omega_i)_e = e_i^*\) (resp. \((\omega_i^{\prime})_e = e_i^*)\) . Then

\[
d \omega_ {k} + \sum_ {i <   j} c _ {i j k} \omega_ {i} \wedge \omega_ {j} = 0 (k = 1, 2, \dots , n)
\]

\[
d \omega_ {k} ^ {\prime} - \sum_ {i <   j} c _ {i j k} \omega_ {i} ^ {\prime} \wedge \omega_ {j} ^ {\prime} = 0 (k = 1, 2, \dots , n).
\]

If \(r < s\) ,

\[
\begin{array}{r l} \left(d \omega_ {k}\right) _ {e} (e _ {r}, e _ {s}) & = - \left(\omega_ {k}\right) _ {e} ([ e _ {r}, e _ {s} ]) \\ & = - \sum_ {i} c _ {r s i} (\omega_ {k}) _ {e} (e _ {i}) \\ & = - c _ {r s k} \\ & = - \sum_ {i <   j} c _ {i j k} (\omega_ {i} \wedge \omega_ {j}) _ {e} (e _ {r}, e _ {s}). \end{array}
\]

The argument is similar for the \(\omega_{k}^{\prime}\) .

